% Style: qhhomework package (~/references/template/sty).
\documentclass[11pt]{article}
\usepackage{qhhomework}

\course{Algebra For Honor Math}
\assignment{W1}
\student{Qinghao Hu}
\duedate{\today}

\begin{document}

\maketitle
\tableofcontents
\newpage

\subsection*{Q3}
\begin{enumerate}[label = (\alph*)]
    \item \begin{equation*}
        \forall m \in \R, \sin^3\paren{2m + 1} = \frac{1}{2}. 
    \end{equation*}
    \item \begin{equation*}
        \exists \theta \in \R, \forall \phi \in \R, \sin \paren{2\theta} < \cos \paren{\phi - \theta}. 
    \end{equation*}
\end{enumerate}

\subsection*{Q4}
\begin{enumerate} [label=(\alph*)]
  \item $$\forall x \in \R, f(x) \leq 3. $$
  \item $$\forall x \in \R, \exists a, b \in \R, x = \frac{a + b }{2}$$. 
\end{enumerate}

\subsection*{Q5}
\begin{enumerate}
  [label=(\alph*)]
  \item This is a true statement. 
  \item This is a false statement. 
  \item This is an open sentence which relies on $q$. 
  \item This is a true statement. 
  \item This is a true statement. 
\end{enumerate}

\subsection*{Q6}
\begin{enumerate}
  [label = (\alph*)]
  \item $S = \mathcal{T}, P(x, y) = P_2(x, y)$. 
  \item $S = \N, P(x, y) = P_3(x, y)$. 
  \item $S = \Z, P(x, y) = P_4(x, y)$. 
  \item $S = \Q, P(x, y) = P_1(x, y)$. 
\end{enumerate}

\end{document}
